\clearpage
\section{Project Domain Analysis}

\subsection{Domain Description}

Electrical power distribution systems are responsible for delivering
electricity from utility providers to consumers while continuously
monitoring energy consumption and ensuring reliable grid operation.
Traditional systems rely on centralized metering infrastructure for
collecting consumption data and managing system operations. The proposed
project introduces a decentralized approach that integrates IoT smart
meters, blockchain technology, and machine learning to provide secure
data management, transparent monitoring, anomaly detection, and
real-time system supervision.

\subsection{Existing System}

In conventional power distribution systems, smart meters or electricity
meters transmit consumption data to a centralized utility database for
storage and processing. System monitoring and energy management are
primarily controlled through centralized infrastructure, creating a
single point of failure and limiting transparency in data management.
Consumers typically have limited access to real-time information
regarding their energy consumption and system status.

\subsection{Proposed System}

The proposed Decentralized Smart Power Grid continuously collects
electricity consumption data from ESP32-based smart meter nodes. The
backend validates incoming telemetry, securely records system
transactions using blockchain technology, and analyzes consumption
patterns through machine learning to detect abnormal usage and possible
power theft. The processed information is stored securely and presented
through a web dashboard, enabling consumers and administrators to
monitor energy consumption, historical records, system performance, and
alerts in real time.